\hwhat{The claim is that regime-III activity coupling, with H38's day-edge drive trimmed, sits in a few strong pairs below percolation. The test infers conditioned pair couplings in 23 regime-III and 23 regime-I units. Each unit has 200 block-shift surrogates, calibrated on synthetic swarms at real schedules. \textbf{Refuted as posed.} Trimming removes 67\% of regime-III bonds (17\% in regime I). The remainder sits at the false-positive floor: 38 bonds against $\approx$23 expected false. The surviving collective excess (10 units, including \#44 and \#51) is a structureless shift of all pairs. The top 10\% of pairs carry a median 11\% of it, and pair structure does not replicate. \#44's fine-tuning team appears in talk, not in activity. Same-lab (provider-outage) pairs carry no excess bonds. At \#51's sampling the method finds strong bonds (recall 0.89). NE43: the edge drive outlives the daily bookend messages. The holdout is not run.}

\hmodel{Inverse Ising (model 01). The spin $s_i(t)=\pm1$ marks agent $i$ active in minute $t$. $\delta_i(b)$ is an agent field per 30-min block. $x_j$ is the block-centered spin; scaffold-unavailable minutes (not started, finished, consolidating, error gap) take the block mean (H38's conditioning). $J_{ij}$ is the equal-time pseudolikelihood bond. A bond is significant when its surrogate $p<1/n_{\rm pairs}$. Percolation uses the Molloy--Reed $\kappa$. CV-C10 is the cross-validated share of the excess covariance $\Delta C$ in the top 10\% of pairs.}

\hmath{\vspace{-8pt}\begin{gather*}
\operatorname{logit}P\big(s_i(t)=1\mid\cdot\big)=2\Big[\delta_i(b)+\textstyle\sum_{j\ne i}J_{ij}\,x_j(t)\Big],\qquad x_j=a_j\,(s_j-\bar s_j^{(b)})\\
z_{ij}=\big(J_{ij}-\langle J_{ij}\rangle_{\rm shift}\big)/\sigma_{ij},\qquad \text{bond if } p(z_{ij})<1/n_{\rm pairs}\\
\kappa=\langle k^2\rangle/\langle k\rangle\ (\text{giant component iff }\kappa>2),\qquad
\mathrm{VR}-1=\textstyle\sum_{i\ne j}\bar C_{ij}\big/\sum_i\bar C_{ii}
\end{gather*}\vspace{-10pt}}

\hobs{../figures/summary_obs.pdf}{(a) Units with a significant edge-trimmed excess (regime III blue, regime I orange). The reference is planted swarms at village sampling (grey $\times$: dense $J=1.2/N$; green $+$: dilute $J=0.5$). Real units sit with the dense cloud: low CV-C10, positive mean bond z. (b) Regime-III significant bonds per unit, raw (grey) and trimmed (blue). Trimming brings them down to the $\approx$1 expected false bond per unit (dotted).}

\htelos{Do not target agent pairs from 1-min activity timing. After the daily start and stop are trimmed, no small clique of agents works in lockstep: 38 bonds stand against $\approx$23 expected false. A weak common mode shared by all pairs remains. It is consistent with a fast shared field as well as with peer coupling. Monitor and steer the common mode, not individual inferred bonds. Look for pairwise structure in who talks in the same minutes: the one known team appears there (p 0.013). The corrected NE43 date is a small data contribution.}

\hbottom{After the operator's schedule is trimmed, regime-III co-activation is a weak, structureless swarm-wide mode, not a dilute ferromagnet; individual bonds sit at the false-positive floor.}

\hresults{{\footnotesize\setlength{\tabcolsep}{2pt}\begin{tabular}{@{}p{0.40\columnwidth}p{0.45\columnwidth}p{0.12\columnwidth}@{}}
\textbf{Prediction} & \textbf{Observed} & \textbf{Verdict}\\\hline
Synthetic: calibrated null; dilute bonds recovered & size 0.054--0.062; recall 0.47 / 0.34 / 0.89 & $\checkmark$ / partial\\
Synthetic: percolation detectable & 0.13 / 0.00 / 0.77 of percolating graphs & $\times$\\
P1 trimmed bonds $\le5\%$ of pairs & 0.95\%; 38 vs $\approx$23 false & $\checkmark$\\
P2 excess concentrated (CV-C10 $\ge0.4$) & 0/10 units; median 0.11 (dense) & $\times$\\
P3 below percolation & 23/23, trivially (power) & $\checkmark$\\
P5 trimming removes bonds in III, not I & 67\% vs 17\% & $\checkmark$\\
NE43 edge drive outlives bookends & $\times0.77$, $\times0.93$ & $\checkmark$\\
\#44 team is a bond cluster & activity p 0.26, 0 bonds; talk p 0.013 & $\times$\\
\#51 bonds recur; same-lab (R2) & r 0.02; OR 0.50 & $\times$\\
NE14 bonds persist II$\to$III & no reliable pairs either side & n/a\\
\end{tabular}}}

\hobsb{../figures/synthetic.png}{Planted swarms at three real schedules (30 replicates each). (a) Raw inference fakes bonds from the day edges; conditioning brings the null to the floor. (b) CV-C10 separates dense from dilute when the excess is significant. (c) Every structure reads as ``below percolation'', even a percolating one. The concentration call is rare for dense coupling and common for dilute coupling.}

\hcaveats{Equal-time inference does not see coupling delayed by minutes, such as read-out-gated replies (H08). The verdict covers equal-time bonds only. Dense $J_0/N$ coupling and a fast shared field look the same here. The post hoc size law favours a field, but it relies on cross-period $N$. Individual bonds need $J\gtrsim0.5$ and $\ge5$ days for detection. CV-C10 became a pooled ratio after the first pass, with no verdict change. All results are rebuilt after the activity-bins join bug; no verdict changed. The holdout is not run.}

\hnext{Fit lagged bonds at the read-out lag. Identify the shared field (human messages, nudges, provider latency) and test the size law at roster joins within \#51. Test dilute against dense in the talk channel, where the \#44 team appears. Run \texttt{confirm.py}.}
